%------------------------------------------------------------------------------%
\documentclass[fleqn,utf8,aspectratio=169,12pt]{beamer}

\usepackage{gaal}

\title{Base e Dimensão}

%------------------------------------------------------------------------------%
\begin{document}

\addtitle

%------------------------------------------------------------------------------%
\section{Base}

%------------------------------------------------------------------------------%
\begin{frame}{Conjunto Gerador}
  Um conjunto
  \[
    G = \left\{ \vec{v}_{1}, \ldots, \vec{v}_{k} \right\}
  \]
  \pause
  é um \emph{conjunto gerador} do \emph{espaço vetorial} $V$

  \pause
  quando todo elemento de $V$ pode ser gerado

  \pause
  como \emph{combinação linear} dos elementos de $G$
  \[
    \pause
    \forall \vec{v}\in V
    \qquad
    \pause
    \vec{v} = \alpha_1\vec{v}_{1} + \cdots + \alpha_k\vec{v}_{k}
  \]
  \pause
  Dizemos que \(G\) gera \(V\)
  \[
    \pause
    V
    = \ger(G)
    \pause
    = \ger\left\{ \vec{v}_{1}, \ldots, \vec{v}_{k} \right\}
  \]
\end{frame}

%------------------------------------------------------------------------------%
\begin{frame}{Base}
  Um conjunto gerador
  \[
    \pause
    \mathcal{B} = \left\{ \vec{v}_{1}, \ldots, \vec{v}_{k} \right\}
  \]
  \pause
  é uma \emph{base} para o espaço vetorial $V$

  \pause
  se seus elementos são \emph{linearmente independentes}
\end{frame}

%------------------------------------------------------------------------------%
\begin{frame}[t]{Base canônica de $\R^2$}
  A \emph{base canônica} para o espaço vetorial \(\R^2\)
  \pause
  é
  \[
    \mathcal{E}
    = \left\{ \vec{i}, \vec{j} \right\}
    \pause
    = \left\{ \vec{e}_{1}, \vec{e}_{2} \right\}
  \]
  \pause
  com
  \[
    \vec{i} = \vec{e}_1 = \vetor{1}{0} \pause \qquad
    \vec{j} = \vec{e}_2 = \vetor{0}{1}
  \]
\end{frame}

%------------------------------------------------------------------------------%
\begin{frame}[t]{Base canônica de $\R^3$}
  A \emph{base canônica} para o espaço vetorial \(\R^3\)
  \pause
  é
  \[
    \mathcal{E}
    = \left\{ \vec{i}, \vec{j}, \vec{k} \right\}
    \pause
    = \left\{ \vec{e}_{1}, \vec{e}_{2}, \vec{e}_{3} \right\}
  \]
  \pause
  com
  \[
    \vec{i} = \vec{e}_1 = \vetortri{1}{0}{0} \pause \qquad
    \vec{j} = \vec{e}_2 = \vetortri{0}{1}{0} \pause \qquad
    \vec{k} = \vec{e}_3 = \vetortri{0}{0}{1}
  \]
\end{frame}

%------------------------------------------------------------------------------%
\begin{frame}[t]{Base canônica de $\R^n$}
  A \emph{base canônica} para o espaço vetorial \(\R^n\)
  \pause
  é
  \[
    \mathcal{E} = \left\{ \vec{e}_{1}, \ldots, \vec{e}_{n} \right\}
  \]
  \pause
  com
  \[
    \vec{e}_1 =
    \begin{pmatrix}[c]
      1      \\
      0      \\
      0      \\
      \vdots \\
      0
    \end{pmatrix}
    \pause
    \qquad
    \vec{e}_2 =
    \begin{pmatrix}[c]
      0      \\
      1      \\
      0      \\
      \vdots \\
      0
    \end{pmatrix}
    \pause
    \qquad
    \vec{e}_3 =
    \begin{pmatrix}[c]
      0      \\
      0      \\
      1      \\
      \vdots \\
      0
    \end{pmatrix}
    \pause
    \qquad
    \cdots
    \qquad
    \vec{e}_n =
    \begin{pmatrix}[c]
      0      \\
      0      \\
      0      \\
      \vdots \\
      1
    \end{pmatrix}
  \]
\end{frame}

%------------------------------------------------------------------------------%
%------------------------------------------------------------------------------%
\begin{question}
  Verifique que
  \[
    \mathcal{E} = \left\{
      \vetor{1}{0},
      \vetor{0}{1}
    \right\}
  \]
  é uma base de $\R^{2}$
\end{question}

%------------------------------------------------------------------------------%
\begin{resolution}{Solução}
  Seja
  \[
    \vec{v} = \vetor{x}{y} \in \R^{2}
  \]
  \pause
  Então
  \[
    \vec{v}
    \pause
    = \vetor{x}{y}
    \pause
    = \vetor{x}{0} + \vetor{0}{y}
    \pause
    = x \vetor{1}{0} + y \vetor{0}{1}
  \]
  \pause
  Logo, $\mathcal{E}$ gera $\R^2$
  \[
    \pause
    \R^2 = \ger\left(\mathcal{E}\right)
  \]
\end{resolution}

%------------------------------------------------------------------------------%
\begin{resolution}{Solução}
  Verificando a independência
  \[
    \pause
    a \vetor{1}{0} + b \vetor{0}{1} = \vetor{0}{0}
  \]
  \pause
  corresponde ao sistema
  \[
    \begin{cases}
      a = 0 \\
      b = 0
    \end{cases}
  \]
  \pause
  Logo $\mathcal{E}$ é LI e gera $\R^{2}$

  \pause
  Portanto $\mathcal{E}$ é uma base
\end{resolution}

%------------------------------------------------------------------------------%
\endinput

%------------------------------------------------------------------------------%
\begin{question}
  Determine se
  \[
    \mathcal{B} = \left\{
      \vetor{1}{2},
      \vetor{3}{1}
    \right\}
  \]
  é uma base de $\R^{2}$
\end{question}

%------------------------------------------------------------------------------%
\begin{resolution}{Solução}
  Primeiro vamos verificar se os vetores são linearmente independentes

  \pause
  Se forem LI verificamos se eles geram $\R^2$
\end{resolution}

%------------------------------------------------------------------------------%
\begin{resolution}{Solução}
  \onslide<+->{Formamos uma matriz com os vetores de $\mathcal{B}$ nas colunas}
  \[
    \onslide<+->{A =
    \begin{pmatrix}
      1 & 3 \\
      2 & 1
    \end{pmatrix}}
  \]
  \begin{align*}
    \onslide<+->{\det(A)}
     \onslide<+->{& = 1\cdot1-3\cdot2 \\}
     \onslide<+->{& = 1-6             \\}
     \onslide<+->{& = -5}
     \onslide<+->{\neq 0}
  \end{align*}
  \onslide<+->{os vetores são LI}
\end{resolution}

%------------------------------------------------------------------------------%
\begin{resolution}{Solução}
  \onslide<+->{Seja
  \(
    \vec{v} = \vetor{x}{y} \in \R^{2}
  \)}

  \onslide<+->{Queremos saber se para quaisquer $x$ e $y$}
  \onslide<+->{existem $a$ e $b$ tais que
  \[
    a\vetor{1}{2} + b\vetor{3}{1} = \vetor{x}{y}
  \]}
  \[
    \onslide<+->{\begin{cases}
    a + 3b = x \\
    2a + b = y
    \end{cases}}
  \]
\end{resolution}

%------------------------------------------------------------------------------%
\begin{resolution}{Solução}
\begin{columns}[t]
\column{0.2\textwidth}
  \begin{align*}
    \onslide<+->{a + 3b & = x      \\}
    \onslide<+->{a      & = x - 3b}
  \end{align*}
\column{0.3\textwidth}
  \begin{align*}
    \onslide<+->{2a + b        & = y              \\}
    \onslide<+->{2(x - 3b) + b & = y              \\}
    \onslide<+->{2x - 6b + b   & = y              \\}
    \onslide<+->{2x - 5b       & = y              \\}
    \onslide<+->{- 5b          & = y - 2x         \\}
    \onslide<+->{b             & = \frac{2x-y}{5}}
  \end{align*}
\column{0.3\textwidth}
  \begin{align*}
    \onslide<+->{a & = x - 3b              \\}
    \onslide<+->{a & = x - 3\frac{2x-y}{5}}
  \end{align*}
\end{columns}
\end{resolution}

%------------------------------------------------------------------------------%
\endinput

%%------------------------------------------------------------------------------%
\begin{question}
  Determine se
  \[
  \mathcal{B}=
  \left\{
  \begin{pmatrix}
  1\\2\\0
  \end{pmatrix}
  ,
  \begin{pmatrix}
  0\\1\\1
  \end{pmatrix}
  ,
  \begin{pmatrix}
  2\\0\\1
  \end{pmatrix}
  \right\}
  \]
  é uma base de $\R^{3}$
\end{question}

%------------------------------------------------------------------------------%
\begin{resolution}{Solução}
  \begin{align*}
  A&=
  \begin{pmatrix}
  1&0&2\\
  2&1&0\\
  0&1&1
  \end{pmatrix}\\
  \det(A)
  &=
  1
  \begin{vmatrix}
  1&0\\
  1&1
  \end{vmatrix}
  +
  2
  \begin{vmatrix}
  2&1\\
  0&1
  \end{vmatrix}\\
  &=1+4\\
  &=5
  \end{align*}
  Como
  \[
  \det(A)\neq0
  \]
  os três vetores são LI.
  Logo
  \[
  \mathcal{B}
  \]
  é uma base de $\R^{3}$
\end{resolution}

%------------------------------------------------------------------------------%
\endinput

%------------------------------------------------------------------------------%
\begin{question}
  Determine se
  \[
    \mathcal{S} =\left\{
      \vetor{1}{0},
      \vetor{0}{1},
      \vetor{1}{1}
    \right\}
  \]
  é uma base de $\R^{2}$
\end{question}

%------------------------------------------------------------------------------%
\begin{resolution}{Solução}
  O conjunto gera $\R^{2}$ pois
  \[
    \vetor{x}{y} = x\vetor{1}{0} + y\vetor{0}{1}
  \]

  \pause
  Porém
  \[
    \vetor{1}{1} = \vetor{1}{0} + \vetor{0}{1}
  \]
  \pause
  Logo os vetores são LD

  \pause
  Portanto $\mathcal{S}$ não é uma base
\end{resolution}

%------------------------------------------------------------------------------%
\endinput


%------------------------------------------------------------------------------%
\section{Dimensão}

%------------------------------------------------------------------------------%
\begin{frame}{Dimensão}
  Se
  \(
    \mathcal{B} = \left\{ \vec{v}_{1}, \ldots, \vec{v}_{n} \right\}
  \)
  é uma base de $V$

  \pause o número de elementos de \(\mathcal{B}\)
  \pause é a \emph{dimensão} e $V$
  \[
    \pause
    \dim(V) = n
  \]
\end{frame}

%------------------------------------------------------------------------------%
\begin{frame}{Dimensão de $\R^{n}$}
  Como a base canônica de $\R^n$ possui $n$ vetores
  \[
    \dim\left(\R^{n}\right) = n
  \]
\end{frame}

%------------------------------------------------------------------------------%
\begin{frame}{Base de $P_{n}$}
  O espaço dos polinômios de grau $n$
  \[
    \pause
    P_{n} = \left\{ a_{0} + a_1x + \cdots + a_nx^n \st a_i \in \R \right\}
  \]
  \pause
  possui a base canônica
  \[
    \pause
    \mathcal{B} = \left\{ 1, x, x^2, \ldots, x^n \right\}
  \]

  \pause
  Como a base possui $n+1$ elementos
  \[
    \pause
    \dim(P_n) = n + 1
  \]
\end{frame}

%------------------------------------------------------------------------------%
\section{Coordenadas}

%------------------------------------------------------------------------------%
\begin{frame}{Coordenadas em uma Base}
  Se
  \(
    \mathcal{B} = \left\{ \vec{v}_{1}, \ldots, \vec{v}_{n} \right\}
  \)
  é uma base de $V$

  \pause
  todo vetor $\vec{v}\in V$ pode ser escrito de maneira \emph{única}
  \[
    \pause
    \vec{v} = a_1\vec{v}_1 + \cdots + a_n\vec{v}_n
  \]

  \pause
  Os escalares
  \[
    \pause
    a_1, \ldots, a_n
  \]
  \pause
  são as \emph{coordenadas} de $\vec{v}$ na base $\mathcal{B}$
\end{frame}

%------------------------------------------------------------------------------%
\begin{frame}{Notação}
  As \emph{coordenadas} de $\vec{v}$ na base $\mathcal{B}$
  \[
    \vec{v} = a_1\vec{v}_1 + \cdots + a_n\vec{v}_n
  \]
  \pause
  são denotadas por
  \[
    \pause
    \left[\,\vec{v}\,\right]_{\mathcal{B}}
    \pause =
    \begin{pmatrix}[c]
      a_{1}  \\
      \vdots \\
      a_{n}
    \end{pmatrix}
  \]
\end{frame}

%------------------------------------------------------------------------------%
%------------------------------------------------------------------------------%
\begin{question}
  Determine as coordenadas do vetor
  \[
    \vec{v} = \vetor{5}{1}
  \]
  na base
  \[
    \mathcal{B} = \left\{
      \vetor{1}{1},
      \vetor{1}{-1}
    \right\}
  \]
\end{question}

%------------------------------------------------------------------------------%
\begin{resolution}{Solução}
\begin{columns}[]
\column{0.4\textwidth}
  \onslide<+->{Procuramos $a$ e $b$ tais que}
  \[
    \onslide<+->{a \vetor{1}{1} + b \vetor{1}{-1} = \vetor{5}{1}}
  \]
  \[
    \onslide<+->{\begin{cases}
      a + b = 5\\
      a - b = 1
    \end{cases}}
  \]

  \onslide<+->{Somando as equações}
  \begin{align*}
    \onslide<+->{2a & = 6 \\}
    \onslide<+->{a  & = 3}
  \end{align*}
\column{0.4\textwidth}
  \onslide<+->{Substituindo na primeira}
  \begin{align*}
    \onslide<+->{a + b & = 5 \\}
    \onslide<+->{3 + b & = 5 \\}
    \onslide<+->{b     & = 2}
  \end{align*}
  \onslide<+->{Logo
  \[
    [\,\vec{v}\,]_{\mathcal{B}} = \vetor{3}{2}
  \]}
\end{columns}
\end{resolution}

%------------------------------------------------------------------------------%
\endinput



%------------------------------------------------------------------------------%
%\section{Lista Mínima}
%
%%------------------------------------------------------------------------------%
%\begin{frame}{Lista Mínima}
%  \begin{enumerate}
%    \item
%  \end{enumerate}
%
%  \vfill
%  Atenção: A prova é baseada no livro, não nas apresentações
%\end{frame}

%------------------------------------------------------------------------------%
\end{document}
%------------------------------------------------------------------------------%
